\usetheme{Berkeley}
\usecolortheme{seahorse}
\usebackgroundtemplate{\includegraphics[width=\paperwidth]{../../paper}}

\beamertemplatenavigationsymbolsempty

\addtobeamertemplate{navigation symbols}{}{
\color{black}\normalsize\textbf\insertframenumber}


\begin{document}

\newcommand{\hrq}[0]{\rule{\textwidth/4}{1pt}}

\begin{titlepage}
\cnam
\begin{center}
{\bfseries

\begin{large}

CONSERVATOIRE NATIONAL DES ARTS ET MÉTIERS

CENTRE RÉGIONAL ASSOCIÉ DE\\
NOUVELLE-AQUITAINE

\hrq

\goal

\vfill
par
\vfill

Marc BENINCA

\hrq

\begin{Large}\fulltitle\end{Large}

Soutenu le 12 Mai 2020

\hrq

JURY

\end{large}
\vspace{1em}
\begin{small}

\begin{tabular}[t]{rllll}
PRÉSIDENT : & Mme & Catherine & BRULATOUT & \textit{\begin{tabular}[t]{@{}l@{}}Maître de conférences associé\\Université de Bordeaux\end{tabular}} \\
MEMBRES : & Mme & Elisabeth & ABRIVAT & \textit{\begin{tabular}[t]{@{}l@{}}Maître de conférences associé\\Université de Bordeaux\end{tabular}} \\
%& ? & Prénom & NOM & \textit{\begin{tabular}[t]{@{}l@{}}Fonction\\Organisme\end{tabular}}
\end{tabular}

\end{small}

}
\end{center}
\end{titlepage}

\cnam

\begin{large}\fulltitle\end{large}

\renewcommand{\contentsname}{Plan}

\cftsetindents{section}{1em}{1.5em}
\cftsetindents{subsection}{2.5em}{2.5em}
\cftsetindents{subsubsection}{5em}{3.25em}
\renewcommand{\cftsecleader}{\cftdotfill{\cftdotsep}}
%\renewcommand{\cftsubsecleader}{\hfill}

\tableofcontents

\pagebreak


\chapter{\ml
{}
{Problématique : maintenance des systèmes d’exploitation}
}

\section{\ml
{}
{Systèmes de fichiers, installés sur partitions, avec accès en écriture}
}

\subsection{\ml
{}
{Système de fichiers conventionnel}
}

\subsection{\ml
{}
{Système de fichiers géré par des recettes configuration}
}

\subsection{\ml
{}
{Système de fichiers avec gestion d’instantanés}
}

\section{\ml
{}
{Images autonomes, sans installation, avec accès en lecture seule}
}

\subsection{\ml
{}
{Amorçage sans gestion de persistance}
}

\subsection{\ml
{}
{Amorçage avec gestion de persistance}
}

\chapter{\ml
{}
{Proposition : fonctionnement autonome incrémental}
}

\section{\ml
{}
{Miroirs de dépôts officiels distribution et éditeurs tiers}
}

\section{\ml
{}
{Construction d’un système de fichiers autonome (Debian GNU/Linux)}
}

\section{\ml
{}
{Encapsulation dans un fichier image}
}

\section{\ml
{}
{Sécurité du fichier image produit}
}

\section{\ml
{}
{Amorçage de fichier(s) image(s) sécurisé(s)}
}

\section{\ml
{}
{Mise à niveau incrémentale}
}

\chapter{\ml
{}
{Automatisations potentiellement implémentables}
}

\section{\ml
{}
{Vérification d’intégrité de dépôts}
}

\section{\ml
{}
{Construction de systèmes de fichiers autonomes complets}
}

\section{\ml
{}
{Création de nouveaux fichiers par mise à jour d’images}
}

\section{\ml
{}
{Génération à la volée de menus de démarrage}
}

\chapter*{}
\addcontentsline{toc}{chapter}{\ml{Summaries}{Résumés}}

\title,\\
\subtitle.

\ml{CNAM Engineer thesis}{Mémoire d’Ingénieur CNAM}, Bordeaux 2021.

\summaries

\thispagestyle{empty}


\end{document}
